\subsection{Executed cover: a planner-latency attack and its fix}\label{sec:cover-study}
Two fixed-weight models, Qwen3:8B and Llama3.1:8B, each run 300 synthetic ticket fixtures per condition: 120 training, 60 validation and 120 test, balanced across six actions. The six conditions yield 3,600 episodes with terminal outcomes, including 1,440 held-out executions. Fixtures and cover draws are paired; request phrasings recur across splits.

\paragraph{Execution and observations}
LangGraph~\cite{langchain2025langgraph} lets each model choose checks, actions, objects, recovery and finish. Seven local MCP Streamable HTTP destinations use TLS~1.3/HTTP/2, a common \path{/mcp} path, and 4,096/16,384-byte request/response padding quanta. A frozen schedule places one or two decoys among early calls. Policies are no cover, proxy discard, read cover, hidden isolated execution, live writes, and visible isolated execution. Proxy discard computes on a disposable copy; hidden isolation retains a clone. Live writes deliberately change the primary fixture as a state-pressure control. Visible isolation forwards clearly marked cover responses to the planner.

The original observer receives destinations and bidirectional ciphertext counts in supplied episode/call windows. LR/RF fit each condition, with validation selecting the primary attacker; position and frozen attacks remain reported. Goal completion checks the requested action/object, prerequisites and target postcondition. Whole-state integrity additionally checks other orders, shared quota and the ledger against the intended-action reference. Each original endpoint retains its twelve-test Holm family. Intervals resample paired fixtures by action, conditional on fitted attacks and the designed workload.

\begin{table*}[t]\centering\small
\caption{Executed cover on the same held-out executions (120 paired fixtures per model and policy; chance 20/120). \emph{Size view}: Byte is the validation-selected size/destination attack; Last is the last-destination rule. \emph{Timing view} (client-receipt call times): Strip removes calls flagged by the label-free gap rule, then applies the frozen no-cover classifier; Jit.\ is the refitted strip attack under 200\,ms half-normal timestamp jitter (mean of five draws); Matched is the best attack after timing-matched decoy dispatch, and +Lat its added latency as a percentage of no-cover wall time. \emph{Execution}: Goal, whole-state integrity (State), completed tasks with state errors, and call-byte ratio. Dashes: not applicable without cover. Supplements S9--S10 give all attacks, intervals and tests.}
\label{tab:cover-main}
\setlength{\tabcolsep}{3.2pt}
\begin{tabularx}{\textwidth}{@{}lYrrrrrrrrrr@{}}\toprule
 & & \multicolumn{2}{c}{Size view} & \multicolumn{4}{c}{Timing view} & \multicolumn{4}{c}{Execution}\\\cmidrule(lr){3-4}\cmidrule(lr){5-8}\cmidrule(l){9-12}
Model & Cover policy & Byte & Last & Strip & Jit. & Matched & +Lat\% & Goal & State & Goal+err & Bytes$\times$ \\\midrule
Qwen3:8B & No cover & 118 & 116 & -- & -- & -- & -- & 114 & 116 & 0 & 1.00 \\
Qwen3:8B & Proxy discard & 62 & 77 & 115 & 104 & 62 & 12.4 & 114 & 116 & 0 & 1.34 \\
Qwen3:8B & Read cover & 62 & 77 & 114 & 103 & 62 & 12.4 & 114 & 115 & 0 & 1.34 \\
Qwen3:8B & Isolated, hidden & 61 & 77 & 114 & 106 & 61 & 12.8 & 114 & 115 & 0 & 1.34 \\
Qwen3:8B & Live write & 62 & 75 & 112 & 102 & 62 & 12.7 & 102 & 7 & 95 & 1.34 \\
Qwen3:8B & Isolated, visible & 63 & 77 & 111 & 98 & 66 & 12.3 & 113 & 114 & 1 & 1.35 \\
\midrule
Llama3.1:8B & No cover & 113 & 79 & -- & -- & -- & -- & 82 & 69 & 13 & 1.00 \\
Llama3.1:8B & Proxy discard & 79 & 67 & 116 & 115 & 79 & 16.4 & 80 & 68 & 13 & 1.29 \\
Llama3.1:8B & Read cover & 72 & 67 & 112 & 110 & 74 & 15.8 & 84 & 71 & 14 & 1.28 \\
Llama3.1:8B & Isolated, hidden & 82 & 67 & 112 & 110 & 82 & 15.7 & 83 & 70 & 13 & 1.28 \\
Llama3.1:8B & Live write & 80 & 65 & 112 & 110 & 80 & 16.4 & 76 & 3 & 73 & 1.29 \\
Llama3.1:8B & Isolated, visible & 81 & 68 & 107 & 105 & 81 & 16.2 & 83 & 73 & 11 & 1.29 \\
\bottomrule\end{tabularx}
\end{table*}


\paragraph{Size-only evaluation credits the cover}
Under the original observer, every cover policy halves recognition. The selected attack scores 61--63/120 for Qwen and 72--82/120 for Llama, against 118 and 113 without cover. By the standard WF-style evaluation, this is a working defense at 1.28--1.35$\times$ call bytes.

\paragraph{A removal attack built on the execution mechanism}
Honest tool calls follow a model decision; the controller dispatches cover directly after a tool response. We therefore add \emph{client-receipt timing}, the start and end time of each call as seen at the Agent's network endpoint. A removal attack flags short-gap calls as cover, deletes them, and classifies the remaining trace with the frozen no-cover classifier. The attacker needs no cover labels. An unsupervised two-component mixture on unlabeled training gaps places the threshold at 0.19\,s (Qwen) and 0.14\,s (Llama). On test it flags all 1,820 decoys and none of the 6,741 honest calls. Stripping restores recognition to 111--115/120 and 107--116/120 (Table~\ref{tab:cover-main}). All ten paired gains survive Holm correction (adjusted $p\leq6.2\times10^{-6}$). A supervised detector trained with cover labels scores within two tasks of the label-free rule (Supplement S9). Each of these analyses reuses the collected executions and was frozen before its own test evaluation (Supplements S9--S10). A generic LR/RF given the same timing features concatenated with the size features does not improve on the size view; none of its ten primary tests rejects. The signal is real, but only an attack that models the execution exploits it, which mirrors the WF experience that attack design, not feature availability, decides a defense's fate~\cite{dyer2012peekaboo,mathews2023critical}.

\paragraph{Why the result is not specific to this testbed}
The leaked quantity is the planner's inference latency, not an implementation constant. Any cover that is generated outside the model loop, at a gateway, proxy, MCP client or tool router, is issued without that latency. To hide, it must reproduce the latency distribution of a particular model, prompt length and load. The margin is large. On validation, cover gaps never exceed 56\,ms, while honest gaps are never shorter than 418\,ms. The attack therefore tolerates network noise. With independent half-normal jitter added to every timestamp, a refitted detector keeps recognition at 107--116/120 up to $\sigma=100$\,ms and 98--115/120 at $\sigma=200$\,ms (Table~\ref{tab:cover-main}, Jit.\ column; Supplement S10). The pre-registered criterion of at least 24 tasks above the size view holds at every tested $\sigma$.

\paragraph{The countermeasure and its price}
If the cause is missing decision time, the fix is to supply it. We simulate \emph{timing-matched dispatch}: each decoy waits for a gap, and lasts for a duration, drawn from the defender's own no-cover history for that model and destination. Because hidden cover never reaches the planner, this delay leaves the planner's behaviour unchanged. A fully refitted attacker then fails. The detector's precision falls to 0.46--0.57, so stripping deletes honest calls and drops to 43--66/120, significantly below the unmatched attack in all ten cells (Holm $p\leq2.1\times10^{-10}$). The best of all reported attacks returns to within three tasks of the size-only score in every cell (Matched column). The price is latency: decoys add 1.1--1.3\,s per covered episode, or 12--16\% of mean no-cover wall time. Timing-matched dispatch therefore turns the size-only operating point back into a real one. The latency cost belongs in the defense budget alongside bytes. This is a receipt-level counterfactual on the collected executions; a live timing-matched cohort and record-level timestamps remain to be measured.

\paragraph{Business outcomes remain a separate constraint}
Proxy, read and hidden-isolated cover retain Qwen state integrity at 115--116/120; adverse live writes reduce it to 7/120. Llama falls from a 69/120 baseline to 3/120. Both live-write integrity losses survive the original correction; completion loss survives only for Qwen. Of 102 Qwen and 76 Llama completed live-write tasks, 95 and 73 have collateral state errors. These counts measure the deliberately adverse control, rather than production cover harm. Similar size-view scores across semantics do not establish equivalence. Together, the two contrasts suggest that schedule and timing determine what the observer learns, while semantics determine what the application pays.

Making isolated responses visible raises mean prompt tokens by 18.1\% and 18.9\%; all original visible-minus-hidden average endpoint tests have adjusted $p=1$. Paths vary between paired runs, so this control remains inconclusive (Supplement S8).
